%% P14 --- Logika, dowod i czytanie twierdzen
%%
%% Kazda liczba w tym programie, ktorej czytelnik nie policzy w pamieci,
%% pochodzi z code/p14_logic_proof.py i trafia na strone przez \val{}.
%% Listing konsolowy to zatwierdzony transkrypt zapisany przez ten skrypt.
%%
%% TEZA: twierdzenie to zalozenia, teza i wiazace je kwantyfikatory, a
%% prawie kazde naduzycie twierdzenia w tej dziedzinie polega na tym, ze
%% czytelnik zatrzymuje teze i gubi cos innego.
%%
%% CO TEN PROGRAM DOSTAJE, odczytane z plikow, a nie zapamietane:
%%   F10  daje "i", "lub", "nie" oraz De Morgana, a jego przyklad z filtrem
%%        jest tutaj kontynuowany -- liczby wypadaja z tabeli prawdy i sa
%%        bramkowane wzgledem tych, ktore tamten program zatwierdzil.
%%   P13  ODSYLA TUTAJ Z NAZWY: dowodzi jednego kierunku twierdzenia o
%%        porzadku topologicznym i mowi, ze czytanie i formulowanie
%%        twierdzen jest uczone w Programie P14.
%%   P05  uzyl ograniczenia sumacyjnego przy swoich liczbach o pojemnosci i
%%        w rigourboxie powiedzial, ze sposob jego uzyskania ma znaczenie.
%%        Sekcja 5 uzywa tego samego ograniczenia do innej wielkosci.
%%   KAZDY rigourbox w ksiazce jest tematem tego programu cwiczonym, zanim
%%        zostal nazwany. NIGDZIE NIE PODAJEMY ICH LICZBY: liczba rosnie z
%%        kazdym napisanym programem, a CLAUDE.md wprost zakazuje podawania
%%        liczby wystapien. Nazwij regule, nie liczbe.
%%
%% CZEGO NIE ROBI, i mowi to na stronie: nie uczy PISANIA dowodow. Ta
%% publicznosc musi czytac twierdzenia, a nie ich dowodzic, i udawanie
%% czegos innego byloby trzecia zla opcja po pominieciu rygoru i po
%% udawaniu go.
%%
%% MNIEJ DO POLICZENIA NIZ W KTORYMKOLWIEK PROGRAMIE, i to fakt o
%% przedmiocie. Co da sie policzyc: kazde zdanie o implikacji i
%% kwantyfikatorach jest rozstrzygniete WYCZERPUJACYM przegladem, ktory
%% jest dowodem, a nie ilustracja. Czego sie nie da, jest tak nazwane, a
%% nie przebrane za pomiar.
%%
%% Blizniak: programs/en/P14-logic-proof.tex. Ramka w ramke, makro w makro,
%% liczba w liczbe; tools/parity.py je porownuje.

\program{Logika, dowód i czytanie twierdzeń}
\label{prog:P14}
\index{proof}
\index{proof!reading a theorem}

\begin{outcomes}
\outcome{1--7}{rozłożyć sformułowane twierdzenie na założenia, tezę i
  wiążące je kwantyfikatory}
\outcome{8--14}{powiedzieć, dlaczego zdanie i jego przeciwstawne to jedno
  twierdzenie, a odwrotne to inne, i co musi pokazać kontrprzykład}
\outcome{15--19}{odczytać kolejność dwóch kwantyfikatorów i powiedzieć, co
  się zmienia po ich zamianie}
\outcome{20--25}{rozpoznać dowód nie wprost i dowód indukcyjny oraz
  powiedzieć, co by każdy z nich złamało}
\outcome{26--31}{powiedzieć o trzech wynikach, które ta dziedzina błędnie
  cytuje, która dokładnie część jest gubiona}
\end{outcomes}

\begin{quiz}
\item \emph{Jeśli to kaczka, to kwacze.} Czy wynika stąd, że coś, co kwacze,
  jest kaczką? \teachesat{9--10}
  \answerto{Nie. To \emph{odwrotne} zdanie i jest innym twierdzeniem; oba
  różnią się w $\val{p14.differ}$ z $\val{p14.rows}$ wierszy. Wynika
  natomiast, że coś, co nie kwacze, nie jest kaczką.}
\item Ograniczenie zachodzi \enquote{z prawdopodobieństwem $\val{p14.conf}$
  procent}. Używasz go $\val{p14.uses}$ razy. Jaka jest szansa, że zachodzą
  wszystkie? \teachesat{28--29}
  \answerto{Około $\val{p14.allhold}$ procent, więc gdzieś zawodzi w
  $\val{p14.anyfails}$ procentach przypadków. Pewność podana raz nie jest
  pewnością wniosku, który budujesz z dwudziestu takich.}
\item \enquote{Dla każdej tolerancji istnieje sieć} oraz \enquote{istnieje
  sieć dla każdej tolerancji}. To samo zdanie? \teachesat{15--16}
  \answerto{Nie, a drugie jest znacznie mocniejsze. Zamiana tych dwóch
  kwantyfikatorów zmienia to, co od czego może zależeć, a twierdzenie o
  uniwersalnej aproksymacji jest pierwszym z tej pary.}
\item Zdanie sprawdzono na czterdziestu kolejnych wejściach i za każdym razem
  zachodzi. Co zostało ustalone? \teachesat{22--23}
  \answerto{To, że zachodzi na tych czterdziestu. Wyrażenie $n^{2}+n+41$ jest
  pierwsze dla każdego $n$ od $0$ do $\val{p14.euler.first}$ wyłącznie, a
  złożone w $\val{p14.euler.first}$, więc czterdzieści potwierdzeń może być
  warte zero.}
\item Twierdzenie o zbieżności zakłada wypukłość funkcji celu. Artykuł cytuje
  je o sieci neuronowej. Co poszło nie tak? \teachesat{26--27}
  \answerto{Zgubiono założenie. Teza jest cytowana, a warunek, pod którym ją
  udowodniono, nie \dash{} a strata sieci neuronowej nie jest wypukła, więc
  twierdzenie po prostu nie ma zastosowania.}
\end{quiz}

% ============================================================
\section{Z czego składa się twierdzenie}
\label{sec:P14-parts}
\index{proof!hypotheses and conclusion}

\begin{fr}
Ten program ma węższy temat, niż się wydaje, i szerszy zysk. Nie nauczy cię
niczego dowodzić. Nauczy cię czytać twierdzenie, które udowodnił ktoś inny, i
mówić dokładnie, co ono orzeka, a czego nie \dash{} a tej umiejętności ta
dziedzina potrzebuje bez przerwy i prawie nigdy jej nie uczy.

Patrzyłeś, jak się to robi. Każdy \emph{rigourbox} w tej książce nazywa coś,
czego program nie robi, i mówi, gdzie robi się to zamiast tego \dash{} zwykle
pierwszy kurs przedmiotu, czasem późniejszy program tutaj, a raz pomiar,
którego nikt nie wykonał. Jeden jest też w tym programie.

Popatrz na ten kształt, a nie na pojedynczy przypadek. Z jakich części składa
się twierdzenie?
\dotline
\nextframe
\end{fr}

\begin{fr}
\begin{ansblock}
Z \textbf{założeń} \dash{} tego, co się przyjmuje. Z \textbf{tezy} \dash{}
tego, co się orzeka. I z \textbf{kwantyfikatorów} \dash{} tego, po czym każde
z nich przebiega.
\end{ansblock}

Trzy części, a prawie każde naduzycie twierdzenia w tej dziedzinie polega na
tym, że jedna z nich ginie w drodze do cytatu.

\begin{note}
Warto od razu wyjaśnić, którą częścią jest dowód. Dowód to argument, że
założenia wymuszają tezę; nie jest jedną z trzech części i można poprawnie
korzystać z twierdzenia, nigdy go nie widząc. To właśnie czyni ten program
możliwym \dash{} i dlatego rigourbox nie jest wykrętem. Zatrzymuje argument, a
daje wszystkie trzy części, czyli wszystko, czego potrzeba, żeby użyć wyniku i
zauważyć, kiedy ktoś inny używa go źle.
\end{note}

\mermaidfig{p14-parts-of-a-theorem}{Trzy części i ta, którą się gubi.}{trzy części}
\end{fr}

\begin{fr}
Rozłóż prawdziwe twierdzenie. Program~\ref{prog:P13} podał takie:

\begin{quote}
Graf skierowany ma porządek topologiczny wtedy i tylko wtedy, gdy jest
acykliczny.
\end{quote}

Nazwij założenia i tezę. W tym zdaniu jest komplikacja, a znalezienie jej to
połowa zadania.
\dotline
\nextframe
\end{fr}

\begin{fr}
\begin{ansblock}
To \emph{dwa} twierdzenia, a każde jest odwrotnym drugiego.
\end{ansblock}

\enquote{Wtedy i tylko wtedy} to skrót na parę zdań podróżujących razem:
\emph{jeśli jest acykliczny, to ma porządek topologiczny} oraz \emph{jeśli ma
porządek topologiczny, to jest acykliczny}. Każde ma własne założenie, własną
tezę i zwykle własny dowód \dash{} a dokładnie to powiedział
Program~\ref{prog:P13}, gdy udowodnił jeden kierunek w jednym zdaniu, a drugi
przekazał tutaj.

\textbf{Założeniem w obu jest \enquote{skierowany}} i słowo to wykonuje
prawdziwą pracę: graf nieskierowany z jedną krawędzią ma marszrutę tam i z
powrotem, więc bez tego słowa twierdzenie jest fałszywe, a cała sekcja 4
Programu~\ref{prog:P13} paruje. Jest to zarazem słowo, które czytelnik
najchętniej pominie, bo stoi przed rzeczownikiem i czyta się jak dekoracja.
\end{fr}

\begin{fr}
O tym wzorcu jest reszta programu, więc nazwijmy go od razu.

\textbf{Najczęstszym nadużyciem twierdzenia jest zatrzymanie tezy i zgubienie
założenia.} Nie niezrozumienie argumentu, nie wątpienie w wynik \dash{}
cytowanie tego, co twierdzenie orzeka, z pominięciem warunków, pod którymi je
udowodniono.

Dlaczego to miałaby być typowa porażka, a nie jakaś inna? Odpowiedz, zanim
pójdziesz dalej; powód dotyczy ludzi, a nie matematyki.
\nextframe
\end{fr}

\begin{fr}
\begin{ansblock}
Bo teza jest tą częścią, o którą komukolwiek chodziło. Założenia to to, co
ktoś musiał przyjąć, żeby ją dostać, a nikt nie cytuje twierdzenia dla jego
założeń.
\end{ansblock}

Teza jest cytowalna, przenośna i krótka. Założenie jest ograniczeniem, jest
nudne i zwykle jest techniczną wtrąconą frazą w środku zdania. Podróżuje więc
teza \dash{} i podróżuje bez tego, co ją uczyniło prawdziwą.

\begin{aibox}
To nie jest hipotetyczna porażka w tej dziedzinie, tylko normalna, a sekcja 5
przerabia trzy jej wystąpienia. Kształt jest zawsze ten sam: wynik
udowodniony pod pewnym warunkiem trafia do wpisu na blogu, na slajd albo do
komentarza w code review bez tego warunku, po czym zostaje zastosowany do
systemu, który go nie spełnia.

Obroną jest nawyk, a nie technika. \textbf{Gdy ktoś cytuje ci wynik, spytaj,
co ten wynik zakładał} \dash{} a jeśli odpowiedzi nie ma w zdaniu, to zdanie
nie jest jeszcze twierdzeniem o twoim systemie.
\end{aibox}
\end{fr}

\begin{fr}
Jednej rzeczy ten program celowo nie robi i mówi to wprost, bo alternatywą
jest zostawić cię w niepewności.

\textbf{Nie uczy pisania dowodów.} Ani w wersji skróconej, ani tych łatwych.
Pisanie dowodów to rzemiosło na kurs i na rok, a prawie nikt wykonujący tę
pracę go nie potrzebuje.

\begin{rigourbox}
Trzy opcje dla książki takiej jak ta były następujące: pominąć rygor
całkowicie, co robi większość praktycznej matematyki dla inżynierów i co
zostawia czytelnika niezdolnego do czytania literatury; uczyć pisania
dowodów, co jest inną książką dla innego czytelnika; albo uczyć
\emph{czytania}, co jest umiejętnością rzeczywiście potrzebną i ułamkiem tej
pracy.

To jest trzecia opcja, a granicę warto nazwać, żebyś wiedział, co będziesz, a
czego nie będziesz miał. Po tym programie umiesz powiedzieć, co twierdzenie
orzeka, co zakłada, co by je obaliło i o czym milczy. Nie umiesz go
udowodnić i nie będzie ci to potrzebne.
\end{rigourbox}
\end{fr}

% ============================================================
\section{Implikacja to nie równoważność}
\label{sec:P14-implication}
\index{proof!implication}

\begin{fr}
Prawie każda para założenie--teza to \emph{implikacja}: jeśli $P$, to $Q$.
Program~\ref{prog:F10} dał \emph{i}, \emph{lub} oraz \emph{nie}; to jest
czwarty spójnik i to jego ludzie rozumieją źle.

Przypadków jest $\val{p14.rows}$, bo $P$ i $Q$ są każde prawdziwe albo
fałszywe. W ilu z nich \enquote{jeśli $P$ to $Q$} jest \emph{fałszywe}?
\dotline
\nextframe
\end{fr}

\begin{fr}
\ans{W jednym}
Tylko wtedy, gdy $P$ zachodzi, a $Q$ zawodzi. Jeśli $P$ nie zachodzi,
implikacja nie orzeka niczego i liczy się jako prawdziwa, co czyta się jak
formalność, a jest całą treścią tego spójnika.

\begin{center}
\begin{tabular}{cccc}
\toprule
$P$ & $Q$ & $P \Rightarrow Q$ & $Q \Rightarrow P$ \\
\midrule
F & F & T & T \\
F & T & T & F \\
T & F & F & T \\
T & T & T & T \\
\bottomrule
\end{tabular}
\end{center}

\textbf{Czyli \enquote{jeśli $P$ to $Q$} to dokładnie zdanie, że nie ma
przypadku, w którym $P$ zachodzi, a $Q$ zawodzi.} Nic o przyczynie, nic o
tym, co się dzieje, gdy $P$ jest fałszywe, nic o tym, że $Q$ jest ciekawe.

Przeczytaj teraz czwartą kolumnę, czyli odwrotne $Q \Rightarrow P$. W ilu
wierszach nie zgadza się z trzecią?
\nextframe
\end{fr}

\begin{fr}
\ans{W $\val{p14.differ}$ z $\val{p14.rows}$}
Są to więc różne zdania i znajomość jednego nie mówi nic o drugim.

\begin{trapbox}
\textbf{Potwierdzanie następnika} to wzięcie $P \Rightarrow Q$,
zaobserwowanie $Q$ i wywnioskowanie $P$. To drugi wiersz tabeli, w którym
implikacja zachodzi, a odwrotne zdanie nie.

W tej dziedzinie przychodzi to zwykle jako diagnoza. \emph{Przeuczenie ciągnie
stratę treningową w dół, a nasza strata treningowa jest w dole, więc się
przeuczamy.} Pierwsze zdanie może być całkowicie prawdziwe, a wniosek i tak
nie wynika: model, który naprawdę nauczył się zadania, też ma niską stratę
treningową.

Test jest mechaniczny i zajmuje sekundę. Zapisz twierdzenie jako \emph{jeśli
$P$ to $Q$}, zapisz, co zaobserwowano, i sprawdź, która z dwóch liter to
jest. Jeśli $Q$, masz zdanie odwrotne i nie masz nic.
\end{trapbox}

Jest błąd będący lustrzanym odbiciem tamtego i to on produkuje fałszywe
uspokojenie zamiast fałszywej diagnozy.

Weź znowu $P \Rightarrow Q$ i przypuść, że obserwujesz, iż $P$ jest
\emph{fałszywe}. Co stąd wynika o $Q$?
\nextframe
\end{fr}

\begin{fr}
\ans{Zupełnie nic}
Wiersze pierwszy i drugi tabeli: przy $P$ fałszywym implikacja zachodzi tak
czy owak, więc $Q$ jest wolne. Wywnioskowanie $\lnot Q$ z $\lnot P$ to
\textbf{zaprzeczanie poprzednikowi} i jest to błąd odwrotnego zdania widziany
z drugiej strony \dash{} istotnie $\lnot P \Rightarrow \lnot Q$ \emph{jest}
zdaniem przeciwstawnym do odwrotnego, więc te dwie pomyłki są jedną pomyłką.

Czyta się to jak bezpieczeństwo, a nie jak twierdzenie, i to pozwala mu
przejść. \emph{Model, który widział zbiór testowy, dostanie podejrzanie dobry
wynik. Ten go nie widział, więc wynik jest wiarygodny.} Drugie zdanie usuwa
jedno wyjaśnienie wysokiego wyniku i zostawia każde inne nietknięte.

Teraz trzecia postać. Czy \enquote{jeśli nie $Q$ to nie $P$} zgadza się z
\enquote{jeśli $P$ to $Q$}?
\dotline
\nextframe
\end{fr}

\begin{fr}
\ans{Tak, we wszystkich $\val{p14.rows}$ wierszach}
To jest zdanie \textbf{przeciwstawne} i nie jest zdaniem pokrewnym, tylko tym
samym zapisanym inaczej \dash{} co skrypt sprawdza, przeglądając każdy wiersz,
więc jest to dowód, a nie ilustracja.

Warto je mieć, bo często jest łatwiejszą połową do wykazania. \emph{Każda
kaczka kwacze} i \emph{nic, co nie kwacze, nie jest kaczką} to jedno zdanie;
jeśli drugie łatwiej sprawdzić, to sprawdzenie go udowodniło pierwsze. Tym
właśnie jest dowód przez kontrapozycję i nie jest to sztuczka.

\mermaidfig{p14-two-rewritings}{Jedno z dwóch przepisań to to samo
zdanie, a drugie to inne.}{dwa przepisania}
\end{fr}

\begin{fr}
Z tabeli wypada jeszcze jedno i jest to najbardziej praktyczna rzecz w tej
sekcji. Czego trzeba, żeby \emph{obalić} \enquote{jeśli $P$ to $Q$}?
\nextframe
\end{fr}

\begin{fr}
\begin{ansblock}
Jednego przypadku, w którym $P$ zachodzi, a $Q$ zawodzi. Dokładnie jednego i
niczego więcej.
\end{ansblock}

Bo zaprzeczeniem $P \Rightarrow Q$ jest $P \wedge \lnot Q$ \dash{} co jest
regułą De Morgana z Programu~\ref{prog:F10} o krok dalej, a skrypt sprawdza to
we wszystkich $\val{p14.rows}$ wierszach.

\textbf{Tym właśnie jest kontrprzykład}, a asymetria zasługuje na uwagę:
wykazanie implikacji może wymagać argumentu o każdym przypadku, a zniszczenie
jej wymaga jednego eksponatu. Dlatego recenzent, który podaje jedno
wystąpienie, zakończył dyskusję, i dlatego \enquote{ale zwykle działa} nie
jest odpowiedzią.

\begin{note}
Program~\ref{prog:F10} zmierzył ten sam spójnik psujący się w filtrze: z
$\val{p14.filter.total}$ rekordów właściwy warunek zachowuje
$\val{p14.filter.right}$, a \code{not spam and not short} zachowuje
$\val{p14.filter.wrong}$, gubiąc $\val{p14.filter.lost}$ bez żadnego objawu.
Skrypt tego programu odbudowuje tamtą populację z tabeli prawdy i przerywa
budowanie, jeśli którakolwiek z liczb nie zgadza się z tą, którą tamten
program zatwierdził \dash{} więc logika tutaj i pomiar tam nie mogą się
rozejść.
\end{note}
\end{fr}

% ============================================================
\section{Kwantyfikatory i kolejność, która zmienia twierdzenie}
\label{sec:P14-quantifiers}
\index{proof!quantifiers}

\begin{fr}
Drugą częścią twierdzenia, którą ludzie gubią, jest kwantyfikator, a
dokładnie jego \emph{kolejność}.

Dwa zdania o relacji między $x$ i $y$:

\begin{itemize}
\item \textbf{Dla każdego $x$ istnieje $y$} takie, że $R(x,y)$.
\item \textbf{Istnieje $y$} takie, że dla każdego $x$ zachodzi $R(x,y)$.
\end{itemize}

Używają tych samych pięciu słów. Czy to to samo twierdzenie? Jeśli nie, to
które jest mocniejsze i dlaczego?
\dotline
\nextframe
\end{fr}

\begin{fr}
\begin{ansblock}
Różne, a drugie jest mocniejsze. W pierwszym $y$ może zależeć od $x$; w
drugim jedno $y$ musi działać dla wszystkich naraz.
\end{ansblock}

Skrypt rozstrzyga to, pokazując relację na $\val{p14.dom}$ elementach \dash{}
każde $x$ w relacji z następnym po okręgu, czyli $\val{p14.rel}$ z
$\val{p14.pairs}$ par \dash{} dla której pierwsze zdanie jest prawdziwe, a
drugie fałszywe. Jeden świadek jest dowodem, że to różne twierdzenia, co jest
asymetrią z sekcji 2 przy pracy.

\begin{note}
Bezpieczny jest tylko jeden kierunek i warto wiedzieć, który.
\textbf{Istnieje-$y$-dla-każdego-$x$ zawsze pociąga
dla-każdego-$x$-istnieje-$y$}, nigdy odwrotnie: świadek uniwersalny służy za
każdego z osobna. Skrypt sprawdza to dla \emph{każdej} relacji na zbiorze
trzyelementowym \dash{} wszystkich $\val{p14.qchecked}$ \dash{} co znów jest
wyczerpującym dowodem, a nie próbką.
\end{note}
\end{fr}

\begin{fr}
Teraz to, co ta dziedzina cytuje najczęściej. Twierdzenie o uniwersalnej
aproksymacji w postaci, w jakiej zwykle się je podaje:

\begin{quote}
Dla każdej funkcji ciągłej $f$ na obszarze ograniczonym i każdej tolerancji
$\varepsilon > 0$ istnieje sieć z jedną warstwą ukrytą, której wyjście jest w
odległości mniejszej niż $\varepsilon$ od $f$ wszędzie na tym obszarze.
\end{quote}

Która to kolejność i co wobec tego może od czego zależeć?
\dotline
\nextframe
\end{fr}

\begin{fr}
\begin{ansblock}
Dla-każdego, a potem istnieje. Sieć może więc zależeć i od $f$, i od
$\varepsilon$ \dash{} i zależy.
\end{ansblock}

Przeczytaj to jeszcze raz z nawykiem z sekcji 2: co to zdanie faktycznie
orzeka? Że dla każdego żądania \emph{coś} je spełnia. A potem przeczytaj to,
czego w zdaniu nie ma wcale.

\begin{itemize}
\item Jak duża musi być ta sieć.
\item Jak ten rozmiar rośnie, gdy $\varepsilon$ maleje albo gdy obszar
  zyskuje wymiary.
\item Jak znaleźć wagi.
\item Czy jakakolwiek procedura uczenia je znajduje.
\item Ile danych by to wymagało.
\end{itemize}

$\val{p14.silent}$ rzeczy, a każda z nich jest tym, o co pytał ten, kto
cytował twierdzenie.

\mermaidfig{p14-exists-is-not-found}{Co obiecuje twierdzenie o istnieniu i
trzy pytania, których nie dotyka.}{obietnica istnienia}
\end{fr}

\begin{fr}
\begin{trapbox}
\textbf{\enquote{Sieci neuronowe potrafią przybliżyć dowolną funkcję} to teza
z usuniętym kwantyfikatorem i usuniętym milczeniem.} Twierdzenie jest
prawdziwe, zdanie jest uczciwym streszczeniem jego tezy, a uzasadnianie nim
wyboru architektury jest błędem, który ten program umie nazwać precyzyjnie:
traktuje twierdzenie o istnieniu jak konstrukcję.

Analogia, która to unaocznia, jest taka, że wielomian dostatecznie wysokiego
stopnia też przybliża dowolną funkcję ciągłą na obszarze ograniczonym i nikt
nie wnioskuje stąd, że wielomiany są dobrą klasą modeli. O sprawie
rozstrzyga wszystko to, o czym twierdzenie milczy.
\end{trapbox}

\begin{warning}
\textbf{To, jak rośnie wymagany rozmiar, jest prawdziwym pytaniem z
prawdziwymi odpowiedziami, a ta książka ani nie przeprowadziła
eksperymentu, ani nie przejrzała tamtej literatury.} Są wyniki w obie strony
\dash{} konstrukcje wydajne dla pewnych klas funkcji i ograniczenia dolne
wykładnicze w wymiarze dla innych.

Powiedzenie, który przypadek stosuje się do konkretnej architektury, wykracza
poza to, co ta książka zrobiła, więc nie jest powiedziane. Ustalony jest tu
kształt zdania, który jest faktem o kwantyfikatorach i nie potrzebuje
eksperymentu.
\end{warning}
\end{fr}

% ============================================================
\section{Trzy kształty i to, czego trzeba, by każdy złamać}
\label{sec:P14-shapes}
\index{proof!induction}
\index{proof!contradiction}

\begin{fr}
Nie musisz pisać dowodów, a musisz rozpoznawać trzy kształty, bo każdy z nich
mówi, jak musiałoby wyglądać obalenie.

Dowód \textbf{wprost} przyjmuje założenia i dochodzi do tezy. Dowód
\textbf{nie wprost} przyjmuje założenia oraz zaprzeczenie tezy i wyprowadza
niemożliwość. Dowód \textbf{indukcyjny} ustala przypadek bazowy i krok od
każdego przypadku do następnego.

Ten drugi to wynik z sekcji 2 przy pracy. Przyjęcie zaprzeczenia
\enquote{jeśli $P$ to $Q$} oznacza przyjęcie czego dokładnie?
\nextframe
\end{fr}

\begin{fr}
\ans{$P$ oraz nie $Q$}
Dlatego dowód nie wprost implikacji zawsze zaczyna się od przypuszczenia, że
założenie zachodzi, a teza zawodzi \dash{} nie jest to wybór stylistyczny,
tylko jedyne, czym zaprzeczenie może być.

Program~\ref{prog:P13} użył takiego dowodu, nie nazywając go: gdyby graf
skierowany miał cykl, każdy wierzchołek na nim musiałby stać po następnym w
koło, a skończona lista tego nie potrafi. Przyjmij założenie, przyjmij, że
teza zawodzi, dojdź do czegoś niemożliwego.
\end{fr}

\begin{fr}
Indukcja jest tu najważniejsza, bo to jej nawyk z inżynierii każe ludziom
sądzić, że ją mają.

Zdanie o każdej liczbie naturalnej. Sprawdzasz je dla $n = 0$, dla $n = 1$ i
dalej, i zachodzi za każdym razem, gdy patrzysz. Dlaczego to nie jest dowód i
co by go dowodem uczyniło?
\dotline
\nextframe
\end{fr}

\begin{fr}
\begin{ansblock}
Bo sprawdziłeś skończenie wiele przypadków, a zdanie jest o nieskończenie
wielu. Dowodem czyni je \emph{krok}: argument, że jeśli zachodzi w $n$, to
zachodzi w $n+1$.
\end{ansblock}

Przypadek bazowy i krok razem sięgają każdego $n$, a żadne z nich osobno nie
sięga żadnego. Taki jest cały kształt, a rozpoznanie go mówi, co atakować:
nieudana indukcja to prawie zawsze krok po cichu zakładający coś więcej, a nie
zły przypadek bazowy.

Oto dlaczego to rozróżnienie nie jest akademickie. Rozważ $n^{2} + n + 41$.

\transcript{p14-forty-confirmations}

\begin{trapbox}
\textbf{Jest pierwsze dla każdego $n$ od $0$ do $\val{p14.euler.first}$
wyłącznie, a złożone w $\val{p14.euler.first}$}, gdzie wynosi
$\val{p14.euler.at}$ \dash{} czyli $\val{p14.euler.factor}$ do kwadratu.

Czterdzieści kolejnych potwierdzeń. Każdy zestaw testów by to przepuścił;
każdy recenzent by uwierzył; a zdanie jest fałszywe.

Uczciwa połowa tego jest warta tyle, co pułapka. W inżynierii tysiąc testów
jest często dokładnie właściwym przyrządem, a ta książka kilka razy użyła
przeglądu jako dowodu \dash{} ale tylko na dziedzinach \emph{skończonych i
wyczerpanych}: cztery wiersze, $\val{p14.qchecked}$ relacji,
$\val{p14.pairs}$ par. \textbf{Sprawdzenie każdego przypadku jest dowodem, a
sprawdzenie wielu przypadków jest przesłanką}, i różnicą nie jest ile, tylko
czy któryś został pominięty.
\end{trapbox}
\end{fr}

\begin{fr}
Każdy kształt psuje się na swój sposób, a wiedza o tym, na który, jest tym, co
czyni rozpoznawanie kształtu cokolwiek wartym. Weź je po kolei: jeśli dowód
jednego z tych trzech rodzajów jest błędny, gdzie najprawdopodobniej leży
błąd?
\dotline
\nextframe
\end{fr}

\begin{fr}
\begin{ansblock}
W dowodzie \emph{wprost} \dash{} w kroku, który nie wynika. W dowodzie
\emph{nie wprost} \dash{} w źle wziętym zaprzeczeniu. W dowodzie
\emph{indukcyjnym} \dash{} w kroku, prawie nigdy w przypadku bazowym.
\end{ansblock}

\begin{itemize}
\item \textbf{Wprost}: któraś linia nie wynika z linii nad nią, a zwykłym
  sprawcą jest krok po cichu korzystający z własności, której nikt nie
  założył.
\item \textbf{Nie wprost}: zaprzeczenie jest złe. Zaprzeczeniem
  \emph{dla każdego $x$ istnieje $y$} jest \emph{istnieje $x$, dla którego
  nie ma $y$}, a pomyłka tutaj sprawia, że cały argument dowodzi czegoś
  innego. Kwantyfikatory z sekcji 3 i $P \wedge \lnot Q$ z sekcji 2 są tym,
  co temu zapobiega.
\item \textbf{Indukcyjny}: krok zachodzi dopiero dla dostatecznie dużych $n$
  albo po cichu potrzebuje zdania dla \emph{wszystkich} mniejszych $n$, gdy
  dostępne jest tylko poprzednie. Przypadek bazowy jest prawie zawsze dobry,
  bo to jego się sprawdza.
\end{itemize}

\textbf{Trzy kształty są więc trzema pytaniami}, a żadne z nich nie wymaga
prześledzenia argumentu w szczegółach. Co zostało zaprzeczone? Czy krok zależy
od tego, że $n$ jest duże? Która linia korzysta z czegoś, czego nie założono?
To właśnie kupuje rozpoznanie kształtu i to całość tego, co ta sekcja
obiecywała.
\end{fr}

% ============================================================
\section{Trzy wyniki, które ta dziedzina cytuje źle}
\label{sec:P14-misquoted}
\index{proof!misquoted results}

\begin{fr}
Wszystko dotąd było maszynerią. Tutaj zostaje użyta, na trzech zdaniach, które
spotkasz.

Pierwsze. Twierdzenie o zbieżności mówi, że spadek gradientu z odpowiednim
krokiem zbiega do minimum globalnego, \emph{o ile funkcja celu jest wypukła}.
Artykuł cytuje je o sieci neuronowej.

Powiedz, co poszło nie tak, i bądź precyzyjny co do tego, której z trzech
części to dotyczy.
\dotline
\nextframe
\end{fr}

\begin{fr}
\begin{ansblock}
Zgubiono założenie. Teza jest cytowana, warunek, pod którym ją udowodniono,
nie jest, a strata sieci neuronowej nie jest wypukła, więc twierdzenie w ogóle
nie ma zastosowania.
\end{ansblock}

Nie \enquote{stosuje się w przybliżeniu} i nie \enquote{stosuje się w
duchu}. Implikacja, której założenie zawodzi, nie orzeka niczego \dash{} to
pierwsze dwa wiersze tabeli z sekcji 2 i to cała ich treść.

\begin{note}
Warto być uczciwym co do tego, co przetrwało, bo \enquote{twierdzenie nie ma
zastosowania} bywa czytane jako \enquote{metoda nie działa}, a to też nie
wynika. Spadek gradientu w widoczny sposób działa na stratach niewypukłych.
Znikła \emph{gwarancja}, a z nią prawo do mówienia \enquote{zbiegnie do
minimum} zamiast \enquote{zbiegał na każdym problemie, który próbowaliśmy}.

To mniejsze twierdzenie i jest uczciwe, a własna reguła tej książki \dash{}
twierdzenie jest zmierzone, przypisane albo oznaczone jako osąd \dash{} to ta
sama dyscyplina zastosowana do zdania zamiast do twierdzenia matematycznego.
\end{note}
\end{fr}

\begin{fr}
Drugie i to ono kosztuje pieniądze. Ograniczenie na generalizację zachodzi
\enquote{z prawdopodobieństwem co najmniej $\val{p14.conf}$ procent po
losowaniu danych treningowych}.

Czytasz artykuł, który stosuje takie ograniczenie $\val{p14.uses}$ razy
\dash{} raz na eksperyment albo raz na model w porównaniu. Przy założeniu, że
losowania są niezależne, jaka jest szansa, że zachodzą wszystkie
$\val{p14.uses}$?
\dotline
\nextframe
\end{fr}

\begin{fr}
\ans{Około $\val{p14.allhold}$ procent}
Szansa, że przynajmniej jedno zawiedzie, wynosi więc $\val{p14.anyfails}$
procent. \textbf{Bardziej prawdopodobne niż nie, że jedno z dwudziestu zdań w
tamtej tabeli jest po prostu fałszywe}, a nic w artykule nie zaznacza które.

Dwie kolejne liczby pokazują ten kształt jasno.

\begin{center}
\begin{tabular}{ll}
\toprule
użyć, zanim staje się to rzutem monetą & $\val{p14.flip}$ \\
pewność na użycie potrzebna do $95\%$ przy $\val{p14.uses}$
  & $\val{p14.needed}$ procent \\
\bottomrule
\end{tabular}
\end{center}

\textbf{Ograniczenie na $\val{p14.conf}$ procent jest rzutem monetą po
$\val{p14.flip}$ użyciach}, a żeby dwadzieścia z nich zachodziło razem, każde
musiałoby zachodzić na $\val{p14.needed}$ procent \dash{} czyli
prawdopodobieństwo porażki $\val{p14.neededdelta}$ procent zamiast $5$, prawie
dwadzieścia razy ciaśniej.

\begin{note}
Założenie niezależności wykonuje wyżej pracę i zwykle jest fałszywe:
dwadzieścia eksperymentów często dzieli dane. Ograniczeniem, które nie
potrzebuje niezależności, jest ograniczenie sumacyjne, którego
Program~\ref{prog:P05} użył przy swoich liczbach o pojemności \dash{} a tutaj
daje ono $\val{p14.uses} \times 5 = \val{p14.union}$ procent, czyli nie mówi
nic.

Nie jest to porażka ograniczenia sumacyjnego; jest to uczciwe stwierdzenie, że
przy dwudziestu użyciach na tej pewności żadna gwarancja nie przeżywa. Oba
rachunki zgadzają się co do tego, co robić, czyli zacieśnić ograniczenie na
pojedyncze użycie zamiast mieć nadzieję.
\end{note}
\end{fr}

\begin{fr}
Trzecim jest uniwersalna aproksymacja, a sekcja 3 już ją rozłożyła, więc
zostaje wzorzec wspólny dla tych trzech.

\begin{aibox}
\textbf{Wypukłość} to zgubione założenie. \textbf{Z wysokim
prawdopodobieństwem} to zgubiony kwantyfikator \dash{} ograniczenie jest
kwantyfikowane po losowaniach danych, a cytowanie go o zbiorze, który masz w
ręku, po cichu zamienia zdanie o rozkładzie na zdanie o próbce.
\textbf{Uniwersalna aproksymacja} to zgubione milczenie: twierdzenie mówi, że
coś istnieje, a cytat sugeruje, że da się to znaleźć.

Po jednym z każdego: założenie, kwantyfikator i przepaść między istnieniem a
konstrukcją. To trzy części z sekcji 1, a czwartego rodzaju pomyłki tu nie
ma, bo nie ma czwartej części.
\end{aibox}

Procedura czytania jest więc dość krótka, żeby użyć jej w pociągu. Co jest
zakładane? Co jest orzekane? Po czym każde z tego przebiega? I o czym zdanie
milczy, a było potrzebne temu, kto je cytował?
\end{fr}

\begin{fr}
Jedna uwaga na koniec i o tym naprawdę był ten program.

Rigourboxy w tej książce nigdy nie były przeprosinami. Każdy z nich nazywa
to, co zostaje przyjęte, przyznaje, że reszta jest gdzie indziej, i mówi gdzie
\dash{} co jest dokładnie tą dyscypliną, którą ten program właśnie uczynił
jawną, ćwiczoną od części Podstawy, zanim została nazwana.

\textbf{Nie brakuje ci dowodu; brakuje ci nawyku pytania, co dany wynik
zakładał.} Ten nawyk kosztuje jedno zdanie, stosuje się tak samo do artykułu,
do kolegi i do strony dokumentacji, i jest jedyną częścią tego przedmiotu,
której będziesz używał co tydzień.
\end{fr}

% ============================================================
\begin{summarybox}
\sumitem{1--2}{\textbf{Twierdzenie to założenia, teza i wiążące je
  kwantyfikatory.} Dowód nie jest żadną z tych trzech części, dlatego z wyniku
  można poprawnie korzystać, nigdy dowodu nie widząc \dash{} i dlatego
  rigourbox nie jest wykrętem.}
\sumitem{3--4}{\enquote{Wtedy i tylko wtedy} to \textbf{dwa} twierdzenia
  podróżujące razem, każde odwrotne do drugiego. W twierdzeniu o porządku
  topologicznym z Programu~\ref{prog:P13} nośnym założeniem jest
  \emph{skierowany} i to słowo czytelnik pomija.}
\sumitem{5--6}{\textbf{Najczęstszym nadużyciem jest zatrzymanie tezy i
  zgubienie założenia}, bo teza jest tą częścią, o którą komukolwiek chodziło.
  Gdy ktoś cytuje ci wynik, spytaj, co ten wynik zakładał.}
\sumitem{7}{Ten program celowo nie uczy pisania dowodów: trzema opcjami było
  pominąć rygor, uczyć pisania albo uczyć \emph{czytania}, a czytanie jest
  umiejętnością potrzebną tej publiczności.}
\sumitem{8--9}{\textbf{$P \Rightarrow Q$ jest fałszywe w dokładnie jednym z
  $\val{p14.rows}$ wierszy} \dash{} $P$ prawdziwe i $Q$ fałszywe. Nie mówi nic
  o przyczynie i nic o przypadku, w którym $P$ zawodzi.}
\sumitem{9--10}{\textbf{Zdanie odwrotne jest innym twierdzeniem}, niezgodnym w
  $\val{p14.differ}$ wierszach. Zaobserwowanie $Q$ i wywnioskowanie $P$ to
  potwierdzanie następnika: \enquote{strata treningowa jest w dole, więc się
  przeuczamy}.}
\sumitem{10--11}{\textbf{A $\lnot P$ nie mówi nic o $Q$.} Wywnioskowanie
  $\lnot Q$ z $\lnot P$ to zaprzeczanie poprzednikowi, czyli błąd odwrotnego
  zdania z drugiej strony \dash{} czyta się jak uspokojenie i to je
  przepuszcza.}
\sumitem{11--12}{\textbf{Zdanie przeciwstawne jest tym samym twierdzeniem}, we
  wszystkich $\val{p14.rows}$ wierszach, i dlatego jego dowód dowodzi
  oryginału.}
\sumitem{13--14}{\textbf{Kontrprzykład to jeden przypadek, w którym założenie
  zachodzi, a teza zawodzi}, bo
  $\lnot(P \Rightarrow Q) = P \wedge \lnot Q$. Wykazanie może wymagać każdego
  przypadku; obalenie wymaga jednego.}
\sumitem{15--16}{\textbf{Kolejność kwantyfikatorów zmienia twierdzenie.}
  Istnieje-$y$-dla-każdego-$x$ jest mocniejsze i pociąga drugie; odwrotnie
  zawodzi, co pokazuje relacja na $\val{p14.dom}$ elementach i potwierdza
  wszystkie $\val{p14.qchecked}$ relacji na trzech.}
\sumitem{17--18}{\textbf{Uniwersalna aproksymacja jest dla-każdego-a-potem-
  istnieje}, więc sieć może zależeć od funkcji i od tolerancji. Milczy o
  $\val{p14.silent}$ rzeczach, w tym o tym, jak duża musi być sieć i czy
  uczenie ją znajduje.}
\sumitem{19}{\enquote{Sieci potrafią przybliżyć dowolną funkcję} traktuje
  twierdzenie o istnieniu jak konstrukcję. Wielomian wysokiego stopnia robi to
  samo i nikt z tego nie argumentuje.}
\sumitem{20--21}{Trzy kształty do rozpoznania: wprost, nie wprost, indukcyjny.
  \textbf{Zaprzeczeniem implikacji jest $P \wedge \lnot Q$}, dlatego dowód nie
  wprost implikacji zawsze tam się zaczyna.}
\sumitem{22--23}{\textbf{Sprawdzenie każdego przypadku jest dowodem;
  sprawdzenie wielu przypadków jest przesłanką.} Wyrażenie $n^{2}+n+41$ jest
  pierwsze dla $\val{p14.euler.first}$ kolejnych wartości i złożone w
  $\val{p14.euler.first}$, gdzie wynosi $\val{p14.euler.factor}$ do kwadratu.}
\sumitem{24--25}{\textbf{Każdy kształt psuje się w swoim miejscu}: dowód
  wprost w kroku korzystającym z czegoś, czego nikt nie założył, dowód nie
  wprost w zaprzeczeniu, a indukcyjny w kroku i prawie nigdy w przypadku
  bazowym. Trzy pytania i żadne nie wymaga prześledzenia argumentu w
  szczegółach.}
\sumitem{26--27}{\textbf{Twierdzenie o zbieżności zakładające wypukłość nie
  mówi nic o sieci neuronowej.} Traci się gwarancję, a nie metodę \dash{} co
  jest własną regułą tej książki o pomiarze zastosowaną do zdania.}
\sumitem{28--29}{\textbf{Ograniczenie na $\val{p14.conf}$ procent użyte
  $\val{p14.uses}$ razy zachodzi wszędzie tylko w $\val{p14.allhold}$
  procentach przypadków}, jest rzutem monetą po $\val{p14.flip}$ użyciach i
  potrzebowałoby $\val{p14.needed}$ procent na każde, żeby dać $95\%$ przy
  dwudziestu.}
\sumitem{30--31}{Trzy błędne cytaty to po jednym z każdej części: zgubione
  założenie, zgubiony kwantyfikator i przepaść między istnieniem a
  konstrukcją. \textbf{Nie ma czwartego rodzaju pomyłki, bo nie ma czwartej
  części.}}
\end{summarybox}

\canyou

\begin{testexercises}
\item \emph{Jeśli model zapamiętuje, jego strata walidacyjna rośnie.} Strata
  walidacyjna wzrosła. Co stąd wynika?
  \answerto{Nic o zapamiętywaniu. To zdanie odwrotne i różni się od oryginału
  w $\val{p14.differ}$ z $\val{p14.rows}$ wierszy \dash{} przesunięcie
  rozkładu albo za wysoki krok uczenia też ją podniosą. Wynika natomiast, ze
  zdania przeciwstawnego, że model, którego strata walidacyjna nie wzrosła,
  nie zapamiętuje \dash{} co jest prawdziwym i użytecznym wnioskiem z tego
  samego zdania.}
\item Zapisz zaprzeczenie zdania \emph{każda partia w tym zbiorze zawiera co
  najmniej jeden przykład pozytywny} i powiedz, co musiałbyś okazać, żeby
  obalić oryginał.
  \answerto{\emph{Pewna partia nie zawiera żadnego przykładu pozytywnego.}
  Zaprzeczenie dla-każdego daje istnieje i odwraca to, co w środku. Żeby
  obalić oryginał, okazujesz jedną taką partię \dash{} nie statystykę o
  partiach, jedną partię.}
\item Ograniczenie zachodzi z prawdopodobieństwem $99\%$ na użycie. Mniej
  więcej ile użyć, zanim szansa, że zachodzą wszystkie, spadnie poniżej
  połowy?
  \answerto{Około $69$: $\num{0.99}^{n} < \num{0.5}$ wymaga
  $n > \ln \num{0.5} / \ln \num{0.99}$. Porównaj $\val{p14.flip}$ użyć przy
  $\val{p14.conf}$ procentach \dash{} zacieśnienie ograniczenia na użycie z
  $5\%$ do $1\%$ kupuje około pięciu razy więcej użyć, co jest wymianą, którą
  autor artykułu musiał zrobić i zwykle jej nie podaje.}
\item \emph{Istnieje krok uczenia działający dla każdej architektury} wobec
  \emph{dla każdej architektury istnieje działający krok uczenia}. Które
  wolałbyś mieć, a które jest prawdopodobne?
  \answerto{Pierwsze jest znacznie mocniejsze i to je domyślna wartość w
  bibliotece orzeka w domyśle; drugie jest niemal puste i to go ktokolwiek by
  bronił. Całą różnicą jest kolejność dwóch kwantyfikatorów, a poradnik
  strojenia, który zsuwa się z drugiego na pierwsze, stawia znacznie większe
  twierdzenie, niż potrafi obronić.}
\item Ktoś mówi, że własność zweryfikowano na dziesięciu tysiącach losowych
  wejść. Co zostało ustalone i kiedy to wystarcza?
  \answerto{To, że zachodzi na tych dziesięciu tysiącach. Wystarcza, gdy
  dziedzina jest skończona i były to wszystkie jej elementy \dash{} tak ten
  program dowodzi swoich zdań o $\val{p14.rows}$ wierszach i
  $\val{p14.qchecked}$ relacjach \dash{} a poza tym jest przesłanką, a nie
  dowodem. Wyrażenie $n^{2}+n+41$ przechodzi czterdzieści kolejnych sprawdzeń
  i jest fałszywe.}
\end{testexercises}

\begin{furtherproblems}
\item Pokaż z tabeli prawdy, że $P \Rightarrow Q$ oraz
  $\lnot P \vee Q$ to to samo zdanie.
  \answerto{Oba są fałszywe dokładnie w wierszu z $P$ prawdziwym i $Q$
  fałszywym, a prawdziwe w pozostałych trzech. Dlatego zaprzeczeniem jest
  $P \wedge \lnot Q$: to De Morgan zastosowany do $\lnot P \vee Q$, czyli
  reguła z Programu~\ref{prog:F10} przy pracy.}
\item Dlaczego \emph{jeśli $P$ to $Q$} liczy się jako prawdziwe zawsze, gdy
  $P$ jest fałszywe, skoro wydaje się, że powinno być nieokreślone?
  \answerto{Bo alternatywa psuje ten spójnik. \emph{Każdy element zbioru
  pustego jest pierwszy} musi być prawdziwe, żeby \emph{jeśli każdy element
  $A$ jest pierwszy i każdy element $B$ jest pierwszy, to każdy element
  $A \cup B$ jest pierwszy} zachodziło we wszystkich przypadkach. Konwencję
  dobrano tak, żeby reguły się składały, i z tego samego powodu iloczyn pusty
  wynosi $1$ w Programie~\ref{prog:F04}.}
\item Założeniami twierdzenia są $A$ i $B$, a tezą $C$. Twój system spełnia
  $A$, ale nie $B$. Co, jeśli cokolwiek, wolno stąd wywnioskować?
  \answerto{Nic z tego twierdzenia, w żadną stronę. Nie mówi ono, że $C$
  zachodzi, i nie mówi, że $C$ zawodzi; założenia zawiodły, więc implikacja
  nie orzeka niczego. Wywnioskowanie, że $C$ zawodzi, to zaprzeczanie
  poprzednikowi, czyli błąd odwrotnego zdania w innym kapeluszu.}
\item Uniwersalną aproksymację formułuje się dla jednej warstwy ukrytej. Co
  musiałaby dodać wersja dla dwóch warstw ukrytych, żeby być mocniejszym
  wynikiem?
  \answerto{Coś o rozmiarze albo o znajdowalności. Czysto jako twierdzenie o
  istnieniu byłaby co najwyżej słabsza \dash{} co potrafi jedna warstwa,
  potrafią dwie \dash{} więc twierdzenie o głębokości jest ciekawe tylko
  wtedy, gdy ogranicza, jak duża musi być sieć albo jak rośnie ten rozmiar. To
  jest pytanie, o którym zdanie o jednej warstwie milczy, i tam jest
  prawdziwa literatura.}
\item Znajdź trzy rigourboxy w programach, które już przeczytałeś. Wskaż
  wspólny im kształt i powiedz, co czytelnikowi wolno zrobić z wynikiem
  podanym w ten sposób.
  \answerto{Każdy podaje tezę, przyznaje, że reszta jest gdzie indziej
  \dash{} dowód, późniejszy program albo raz pomiar, którego nikt nie wykonał
  \dash{} i nazywa gdzie. Czytelnikowi wolno użyć wyniku tak, jakby widział
  dowód \dash{} wszystkie trzy części są obecne \dash{} a nie wolno go
  rozszerzać, osłabiać założenia ani stosować poza jego kwantyfikatorami, bo
  to są ruchy, które argument musiałby uzasadnić.}
\end{furtherproblems}
